\section{Introduction}
\label{sec:intro}
\PARstart{S}{ervice} continuity concerns whether a person can continue the task they need during a disruption. A device may lose its bearer yet retain a cached lesson, document editor, or local development tool. Conversely, a connected bearer alone does not establish that an application meets its deadline, has the needed data, or fits the available compute and memory. IMT-2030 identifies service-oriented resilience and ubiquitous connectivity as important directions~\cite{itur_m2160}; versioned 5G service requirements already distinguish performance expectations among services~\cite{3gpp_ts22261}. These frameworks motivate an application-sensitive evaluation, but they do not validate the thresholds or priors used in a particular simulation.

For resource-constrained edge devices, a useful evaluation must preserve both network and workload state. This article asks a deliberately bounded question: \emph{what does the currently implemented continuity benchmark establish about declared local workload support and simulated useful-service occupancy, and which apparently stronger conclusions are unsupported?} The evidence is digital: a frozen Device OS service-profile corpus and a separate seeded generator in the 7GC Digital Twin repository. The term digital twin identifies the repository and its research trajectory; the evaluated generator has no live physical counterpart or calibrated radio model.

The study makes three contributions. First, it formalizes and reproduces a four-level classification over 21 stored profile snapshots, including an explicit unsupported workload and the distinction between a scenario name and the state actually exercised. Second, it documents a 30-seed generator at the level of its random draws, synthetic class priors, and statistical unit, allowing exact reproduction and analytic interpretation of its two primary outcomes. Third, it audits the dependency structure and meaning of secondary metrics, identifies an outage-entry counter mislabeled as recovery, and states a prospective path to physical validation. These contributions concern benchmark semantics and reproducibility; they establish neither a new optimal controller nor a performance improvement over published systems.

The Gary flagship scenario supplies named disruption contexts such as neighborhood outages and congested school connectivity. It is an authored research scenario, rather than an operating community network. Five context names map to stored bearer labels. That mapping supplies a transparent scenario vocabulary, but does not model storm propagation, neighborhood infrastructure, payment transaction integrity, or RF interference. The analysis retains unfavorable and inconclusive findings instead of treating successful execution as evidence of general resilience.

The remainder distinguishes prior work, the two evidence paths, experimental versioning, reproduced outcomes, and threats to validity. A separate author-review supplement records exact hashes, commands, a data dictionary, citation checks, and claims requiring future evidence.
